\chapter{Introduction}
\label{ch:intro}

\section{Purpose and scope}

A golf launch monitor answers two questions for every shot: \emph{what did
the club do through impact} (delivery), and \emph{what did the ball do after
impact} (launch and flight). Commercial systems answer these questions with
strikingly different sensor architectures --- continuous-wave Doppler radar,
multi-camera photogrammetry, and, increasingly, fusions of the two --- yet
they all implement the same underlying physics: an oblique-impact collision
model connecting club delivery to ball launch, and an aerodynamic model
connecting ball launch to flight.

This review is written to support the development of
\textbf{OpenFlight}\footnote{\url{https://github.com/jewbetcha/openflight}
--- AGPL-3.0; a DIY launch monitor built on the OmniPreSense OPS243-A
24\,GHz Doppler radar, RFbeam K-LD7 angle radars, a sound-trigger board, and
a Raspberry~Pi~5.}, an open-source launch monitor. Its goals are:

\begin{enumerate}[itemsep=2pt]
  \item catalogue \emph{how each major commercial system works}, at the level
        of radar physics, imaging geometry, and signal processing;
  \item establish, parameter by parameter, \emph{what is measured directly,
        what is inferred through a model, and what is essentially
        estimated} --- especially for club-delivery data such as face angle
        and club path;
  \item survey the governing \emph{patent landscape}, since the strongest
        public documentation of these proprietary methods is found in the
        patents themselves, and since freedom-to-operate matters to an
        open-source project;
  \item collect the \emph{physics and algorithms} (D-plane, impact mechanics,
        gear effect, aerodynamic coefficients, spectral spin estimation,
        stereo triangulation, Kalman filtering) that any implementation
        needs; and
  \item translate all of the above into \emph{concrete design guidance} for
        OpenFlight's radar-first architecture and its possible optical
        extensions.
\end{enumerate}

\section{The two sensing families}

\textbf{Doppler radar} systems (TrackMan, FlightScope, Garmin Approach R10,
Full Swing KIT) illuminate the hitting area and downrange volume with a
microwave carrier and extract target radial velocity from the Doppler shift.
Multiple receive antennas turn a velocity sensor into a 3D tracker via phase
interferometry (\cref{ch:radar}). Radar excels outdoors: it tracks the
entire flight, so carry and curvature are \emph{observed}, not modeled. Its
weaknesses are at the club face --- radar cannot see face orientation
directly --- and indoors, where a screen truncates the observable flight.

\textbf{Photometric (camera)} systems (Foresight GC-series, Uneekor,
SkyTrak, ProTee VX, Garmin R50) capture a burst of high-speed, IR-strobed
stereo images over the first $\sim$30\,cm of ball flight and reconstruct
position and orientation photogrammetrically (\cref{ch:camera}). They
measure launch conditions --- including 3D spin from dimple-pattern rotation
--- essentially perfectly for simulator purposes, then \emph{model} the
flight. Their weakness is the mirror image of radar's: everything downrange
of the capture volume is simulated, and club data requires either fiducial
stickers on the face or an overhead viewing geometry.

The market's convergent evolution is the single most instructive fact for a
new design: \emph{every} high-end vendor has concluded that neither modality
suffices alone. TrackMan added cameras to its radar (OERT, \cref{sec:oert});
FlightScope added Fusion Tracking cameras; SkyTrak added radar to its
camera; Rapsodo pairs radar with impact cameras; Full Swing feeds a camera
into its radar ML pipeline. \Cref{ch:survey} tabulates the landscape.

\section{Reading guide}

\Cref{ch:params} fixes terminology and coordinate conventions
(TrackMan's definitions, the de~facto industry standard).
\Cref{ch:impact} develops the impact physics connecting club delivery to
ball launch --- the D-plane model, spin generation, and gear effect --- which
every monitor uses either forward (simulation) or inverse (parameter
estimation). \Cref{ch:radar,ch:camera} treat the two sensing families in
depth. \Cref{ch:survey} surveys the commercial devices.
\Cref{ch:patents} maps the patent landscape with expiry status.
\Cref{ch:flight} covers ball-flight aerodynamics and trajectory estimation.
\Cref{ch:accuracy} reviews the independent validation literature.
\Cref{ch:implications} distills the implications for OpenFlight.

\begin{keypoint}
Throughout, we use a right-handed coordinate system for a right-handed
golfer: $x$ down the target line, $y$ up, $z$ to the golfer's right.
Angles are positive right/up. All club-delivery values are referenced to
the moment of \emph{maximum ball compression}; all ball-launch values to
the moment of \emph{separation from the face}.
\end{keypoint}
